\documentclass[11pt,a4paper]{article}
\usepackage[margin=25mm]{geometry}
\usepackage{amsmath,amssymb,amsthm,booktabs}
\newtheorem{lemma}{Lemma}\newtheorem{theorem}{Theorem}
\title{A common involution and the pair count of type $(7,8)$}
\author{Human Compression project}\date{}
\begin{document}\maketitle
Write $G=O'N$, and let $x_J(p,q)$ count ordered pairs $(r,s)$ in a
subgroup $J$ with $|r|=p$, $|s|=q$ and $|rs|=2$. Write $e_J(p,q)$
for the number which generate $J$. We prove the fixed Connor--Leemans
Table~1 entry
\[                         e_G(7,8)/(2|G|)=1848.                 \tag{1}\]
The complete subgroup roster, the primary ATLAS class data and the
previously acquired character evaluations are borrowed inputs.
All row numbers below refer to the fixed subgroup roster.

\begin{lemma}\label{atmost}
A pair of type $(7,8)$ lies in at most one maximal subgroup of row 7.
\end{lemma}
\begin{proof}
The row 7 maximal subgroups are the involution centralizers. For
$|r|=7$, the primary class data give $|C_G(r)|=1372=4\cdot343$ or
$49$. In the first case the same power data show that $r$ is a power
of an element of order 28. Its centralizer therefore has a cyclic
Sylow 2-subgroup of order four. In the second case it has odd order.
For $|s|=8$, both possible centralizers have order 32. Consequently
$C_G(r,s)$ is a 2-group contained in a cyclic Sylow 2-subgroup of
$C_G(r)$, and is cyclic of order at most four. It contains at most
one involution. A row 7 maximal contains both entries exactly when
its defining involution belongs to $C_G(r,s)$, proving the assertion.
\end{proof}

\begin{lemma}\label{alignment}
Each subgroup in rows 91, 92, 154 and 155 has a unique maximal
overgroup in $G$.
\end{lemma}
\begin{proof}
The complete ancestor lists show that every maximal overgroup of
one of these subgroups is reached through a parabolic subgroup $P$
in row 9 or 10. Such a $P$ has order 32928 and a normal elementary
abelian subgroup $V$ of order 49, with $|P/V|=672$.
We first make this structural input explicit. In the three-dimensional
projective model, each of the 57 point stabilizers and 57 plane
stabilizers in $L_3(7)$ has order 32928. These 114 subgroups are distinct:
a point stabilizer fixes just its specified point, and its full
$\mathrm{GL}_2(7)$ action on the quotient preserves no plane through
that point; its unipotent radical excludes an invariant complementary
plane. Duality gives the corresponding plane assertions. The roster
lists precisely 114 maximal subgroups of this order in each fixed
$L_3(7)$, all in the appropriate row 9 or 10. Thus these are exactly
the listed parabolics. In coordinates for a point or plane, the two
free entries in the unipotent radical give $V\cong C_7^2$ and
$|P/V|=672$.

Let $L$ be one of the four subgroups in the statement. Its order is
784 or 392, so its 7-part is 49. Since the 7-part of $|P/V|$ is
seven, $|L\cap V|$ is at least seven, and this intersection is normal
in $L$. If it had order seven, $L$ would normalize a cyclic subgroup
of that order. There are exactly two such classes in $G$. Their
normalizers have orders 294 and 8232. The former has no element of
order eight. The latter is the accepted flag normalizer in row 22;
the complete literal model has no element of order eight either.
Every $L$ under consideration contains an element of order eight,
as its complete cyclic ancestor join records. This rules out an
intersection of order seven.

It follows that $V\leq L$. As a normal Sylow 7-subgroup of $L$, it
is unique and characteristic. Repeating the argument for any other
parabolic $P'$ containing $L$ shows that its radical is the same
$V$. The three subgroup classes of order 49 have normalizer orders
32928, 2058 and 32928. The middle one is impossible because its
normalizer would contain $L$ and hence an element of order eight.
Therefore $|N_G(V)|=32928$, and $P=P'=N_G(V)$.
Finally each row 9 or 10 subgroup has one actual simple overgroup
in row 4 or 5, and each of those has one actual maximal overgroup
in row 2 or 3. The direct class-length incidences prove these
uniqueness assertions. Thus every maximal overgroup of $L$ is the
same subgroup.
\end{proof}

\begin{lemma}\label{unique}
A proper subgroup generated by a pair of type $(7,8)$, and contained
in no row 7 maximal, has a unique maximal overgroup.
\end{lemma}
\begin{proof}
The cyclic order seven classes are rows 569 and 570; the cyclic
order eight classes are rows 559 and 560. Intersect their complete
ancestor sets, remove $G$ and the descendants of row 7. This leaves
\[
2,3,4,5,9,10,13,14,21,45,46,53,54,91,92,93,94,95,96,154,155,172,174.
\]
The character evaluation gives $x_M(7,8)=0$ for row 21, so this
class cannot contain the pair. Rows 2 and 3 are maximal.
Apart from the four rows in Lemma~\ref{alignment}, all remaining
rows have the following unique immediate-parent classes:
\[
\begin{array}{c|rrrrrrrr}
\text{child}&4&5&9&10&13&14&45&46\\
\text{parent}&2&3&4&5&9&10&13&14\\\hline
\text{child}&53&54&93,94&95,96&172&174\\
\text{parent}&9&10&53&54&4&5.
\end{array}
\]
For every displayed edge, the parent class length times its number
of child subgroups equals the child class length. Each actual child
therefore has just one parent. Composing these proved uniqueness
statements, together with Lemma~\ref{alignment}, proves the claim.
No multiplicity is obtained by multiplying arbitrary cover paths.
\end{proof}

\begin{lemma}\label{inside}
Suppose a pair of type $(7,8)$ lies in a projective maximal $M$ and
in a row 7 maximal $C_G(t)$. Then $t\in M$.
\end{lemma}
\begin{proof}
Put $L=\langle r,s\rangle$. The intersection of the two cyclic
ancestor sets with the descendants of row 7 and rows 2 or 3 consists
of exactly
\[
65,66,110,111,112,113,115,116,117,118,171,173,175,176.             \tag{2}
\]
We verify the required assertion for every possible class of $L$.
The block group $H$ of order 1344 has a perfect subgroup
$D\cong\mathrm{SL}_2(7)$ of order 336. In the explicit block model,
$D$ is the image of the determinant-one block; the scalar quotient
of order three intersects it trivially. Perfectness follows directly from the elementary generators: conjugation
by $\operatorname{diag}(2,2^{-1})$ makes every upper and lower
unipotent a commutator, and these generate $\mathrm{SL}_2(7)$.
Its central involution $z$
is central in $H$, including the graph automorphism. Every subgroup
of $H$ of index at most four contains $D$: the action of $D$ on
its cosets has image in the solvable group $S_4$, and a perfect
image in a solvable group is trivial. Thus each such subgroup
contains $z$. The classes 65, 66, 112, 113, 115--118, 175 and 176
have an actual block overgroup and index at most four in it.
For any $L$ in these classes, its central involution $z$ must be
$t$ by Lemma~\ref{atmost}. Hence $t\in L\leq M$.

The outer-involution centralizer in a projective maximal is
$K\cong\mathrm{PGL}_2(7)\times C_2$. It belongs to row 110 or 111.
Here is an explicit identification from the maximal roster.
The only larger maximal subgroup orders divisible by 672 are
1876896 and 1344. The former is the socle, which cannot contain the
outer involution in $K$. If $K$ lay in a block group of order 1344,
its index two would force $D\leq K$. This is impossible because
$D$ is perfect of order 336 whereas $K'$ has order 168.
Thus $K$ is maximal, and the only appropriate maximal row of order
672 is 110 or 111. Consequently all members of these two classes
have the central involution of their direct factor. When $L$ is
one of them, this involution equals $t$ and belongs to $M$.

The remaining classes are 171 and 173. Each has one actual parent
$K$ in row 110 or 111, of index two, and one actual simple parent
$S$ in row 4 or 5. The direct incidence equalities prove both
uniqueness assertions. Since $L\lhd K$ and $|N_G(L)|=672=|K|$,
we have $K=N_G(L)$. Uniqueness of $S$ implies that $K$ normalizes
$S$. Its normalizer has order 3753792 and is the unique maximal
parent of $S$. Thus that maximal contains $K$.
Every overgroup path from $L$ starts through $K$ or $S$; in either
case its projective maximal contains $K$ and hence the central
involution of $K$. That involution centralizes $L$, so it is $t$.
This proves the assertion in every class of (2).
\end{proof}

\begin{theorem}
Let $H$ and $K$ be the block and outer-involution centralizers above.
Then
\begin{align*}
 e_G(7,8)={}&x_G(7,8)-\sum_M[G:M]x_M(7,8)\\
 &+2\cdot342868680\bigl(x_H(7,8)+2x_K(7,8)\bigr),             \tag{3}
\end{align*}
where $M$ ranges over the thirteen maximal conjugacy classes.
The resulting ordered automorphism-orbit count is 1848.
\end{theorem}
\begin{proof}
A generating pair belongs to no maximal subgroup. A nongenerating
pair with no row 7 overgroup belongs to exactly one maximal subgroup,
by Lemma~\ref{unique}. A pair with a row 7 overgroup belongs to
exactly one of that class, by Lemma~\ref{atmost}; its excess
multiplicity in the maximal sum is exactly the number of its other
maximal overgroups.

The only maximal classes with nonzero $x_M(7,8)$ are rows 2, 3 and 7.
Within either projective maximal $M\cong L_3(7):2$, Lemma~\ref{inside}
shows that the pairs having a row 7 overgroup are counted by summing
$x_{C_M(t)}(7,8)$ over all involutions $t\in M$. There is no repeated
pair in this sum by Lemma~\ref{atmost}. The accepted centralizer
calculation gives 2793 inner involutions with centralizer $H$ and
5586 outer involutions with centralizer $K$. Each projective maximal
class has 122760 members. Their product
$122760\cdot2793=342868680$ gives (3).

The literal calculation gives $x_H(7,8)=672$ and $x_K(7,8)=1344$.
The character formula gives
\[
x_G(7,8)=1734970379788800,\qquad
\sum_M[G:M]x_M(7,8)=34100347438080.
\]
Substitution into (3), followed by division by $2|G|=921631011840$,
gives 1848. This division counts automorphism orbits of ordered
generating pairs, since the automorphism action on them is free.
No quotient by interchange of the entries is taken.
\end{proof}

\paragraph{The finite and source certificates.}
The retained literal evaluation reconstructs both small groups from
their generators, partitions every element of order eight into
conjugacy classes, and tests every partner of each representative.
The independent reader uses pointwise permutation composition and
compares all returned classes and partners. It also checks the
complete flag normalizer spectrum, which excludes order eight in
the order-8232 normalizer. All ancestor, descendant, normalizer and
direct-incidence assertions above are joined to the exact native
roster pages; no arbitrary cover-chain product is used.

The numerical result is due to Connor and Leemans. The present
contribution is the common-involution argument, the normal-Sylow
alignment and the exact maximal-subgroup correction. The subgroup
roster, primary centralizer data, classical projective models and
accepted character evaluations remain borrowed premises. This
settles one additional numerical entry, and supplies no chirality,
whole-source, open-question or Lean credit.
\end{document}
